% Standalone IEEE-compatible replacement table and text.
% New 85-point case: do not append these values to the former 50-point row.
\begin{table}[t]
\caption{Andermatt local DH routing: generated-to-installed ratios on a new 85-transfer-point case. Unity indicates agreement.}
\label{tab:andermatt_topotherm_accuracy}
\centering
\footnotesize
\begin{tabular}{lrr}
\hline
Measure & Urban4 & Topotherm \\
\hline
Route length / area & 0.880 & 0.622 \\
Route length / connection & 0.880 & 0.622 \\
Connection spacing (fixed input) & 1.000 & 1.000 \\
Building-to-route distance & 1.390 & 1.750 \\
Mean branch length & 2.412 & 1.760 \\
Terminal share & 0.975 & 1.004 \\
\hline
\end{tabular}

\vspace{2pt}
\parbox{\columnwidth}{\scriptsize Both methods receive identical observed connections, inlet, scenario loads and OSM road candidates. The first two measures are normalisations of the same route extent; spacing is an input control. Topology descriptors use a common 0.5-m geometry protocol and degree-2 suppression.}
\end{table}

We evaluated conditional district-heating route reconstruction in Andermatt against active cadastral routes, using 85 observed transfer points, one local inlet and identical OSM road candidates for Urban4 and Topotherm. Uniform 50-kW scenario loads were imposed at all points. Urban4 and Topotherm generated 9.56 and 6.75~km of routes against 10.86~km mapped in the common window. At a 25-m tolerance, Urban4 covered 68.1\% of installed route length and Topotherm 62.7\%; the reverse fractions were 74.4\% and 86.7\%. Mean generated-to-installed route distances were 28.0 and 15.0~m. Urban4 better reproduced mapped extent and building proximity, whereas Topotherm better localised its more compact network. Topotherm reached a 1\% MIP gap in 422~s. These results quantify a routing tradeoff conditional on supplied connections and scenario loads; they do not validate customer selection, measured heat demand or hydraulic performance.

% Cite the Lisag AG / Geoinformation Kanton Uri "Fernwaerme Leitungen"
% WFS route dataset (retrieved 2026-11-03), and the official Topotherm
% software/paper. The SFOE network inventory alone is not a detailed route source.
